% Part: sets-functions-relations
% Chapter: size-of-sets-complete

\documentclass[../../../include/open-logic-chapter]{subfiles}

\begin{document}

\olchapter{sfr}{siz}{O Tamanho dos Conjuntos}

\begin{editorial}
Este capítulo discute enumerações e conjuntos contáveis e incontáveis.
Várias seções vêm em duas versões: uma mais elementar, que
toma enumerações como listas, ou sobrejeções de $\PosInt$; e uma
mais abstrata que define enumerações como bijeções com $\Nat$.
\end{editorial}

\olimport{introduction}

\olimport{enumerability}

\olimport{zig-zag}

\olimport{pairing}

\olimport{pairing-alt}

\olimport{non-enumerability}

\olimport{reduction}

\olimport{equinumerous-sets}

\olimport{comparing-size}

\olimport{schroder-bernstein}

\begin{editorial}
As seções \olref[sfr][siz][enm-alt]{sec},
\olref[sfr][siz][nen-alt]{sec}, \olref[sfr][siz][red-alt]{sec} são
versões alternativas da \olref[sfr][siz][enm]{sec},
\olref[sfr][siz][nen]{sec}, \olref[sfr][siz][red]{sec} de Tim
Button para uso em seu texto Open Set Theory. São um pouco mais
avançadas e usam uma definição diferente de enumerabilidade mais
adequada em um contexto de teoria dos conjuntos (isto é, bijeção com $\Nat$ ou um
segmento inicial, em vez de ser listável ou ser a imagem de uma
função sobrejetora de $\PosInt$).
\end{editorial}

\olimport{enumerability-alt}
\olimport{non-enumerability-alt}
\olimport{reduction-alt}

\OLEndChapterHook

\end{document}
